% 矩阵方程（二）：Lyapunov方程X-A^TXA=Q与A^TX+XA=-Q（离散累计与唯一可解条件、一次Schur变换下的对称递推与检验、级数解、半正定右端的正定条件与谱证书）

\section{Lyapunov方程}
\label{sec:ls-lyapunov}

第~\ref{chap:linear-recurrence-stability}章研究了反复作用$\bm x_{k+1}=\bm A\bm x_k$的长期行为。若每一步以二次量$\bm x_k^\top\bm Q\bm x_k$计量状态的大小，能否用一个矩阵$\bm X$直接写出全部步骤的累计量$\bm x_0^\top\bm X\bm x_0$？本节由这一离散累计引出Lyapunov方程，并把相应的连续型方程作为Sylvester方程的特例处理。

全节统一采用两种写法：离散型$\bm X-\bm A^\top\bm X\bm A=\bm Q$，连续型$\bm A^\top\bm X+\bm X\bm A=-\bm Q$，其中$\bm Q$对称，未知量$\bm X$与上一节一样记为$\bm X$。两者都是左乘$\bm A^\top$、右乘$\bm A$，所以只需对$\bm A$作一次Schur分解，回代次序也相同。离散型的正定性论证只用矩阵幂和级数；连续型正定解的积分证明需要矩阵指数，本节只陈述条件并给出代数检验。

\subsection{离散型与连续型方程}
\label{sec:ls-lyapunov-solvability}

设$\bm A\in\mathbb R^{d\times d}$，$\bm Q=\bm Q^\top\succeq0$，$\bm x_k=\bm A^k\bm x_0$。前$N+1$步的累计量为
\[
 \sum_{k=0}^{N}\bm x_k^\top\bm Q\bm x_k=\bm x_0^\top\bm X_N\bm x_0,\qquad
 \bm X_N=\sum_{k=0}^{N}(\bm A^\top)^k\bm Q\bm A^k ,
\]
其中$\bm X_N$与初值无关。把和式的首项单独写出，其余各项提出左右的$\bm A^\top$与$\bm A$，得递推$\bm X_{N+1}=\bm Q+\bm A^\top\bm X_N\bm A$。若$\bm X_N$收敛到$\bm X$，对递推取极限，极限应满足离散型\term{Lyapunov方程}（Lyapunov equation）
\begin{equation}
 \bm X-\bm A^\top\bm X\bm A=\bm Q .
 \label{eq:ls-discrete-lyapunov}
\end{equation}
它是矩阵空间上固定更新$\bm X\mapsto\bm Q+\bm A^\top\bm X\bm A$的不动点方程。

方程~\eqref{eq:ls-discrete-lyapunov}也可以逐步读出。若$\bm X$是解，则沿轨迹有
\[
 \bm x_{k+1}^\top\bm X\bm x_{k+1}-\bm x_k^\top\bm X\bm x_k
 =\bm x_k^\top(\bm A^\top\bm X\bm A-\bm X)\bm x_k
 =-\bm x_k^\top\bm Q\bm x_k ,
\]
即二次量$\bm x^\top\bm X\bm x$每走一步恰好减少这一步的计量$\bm x_k^\top\bm Q\bm x_k$。把这些等式从$k=0$加到$N$，就回到上面的累计量。

离散型方程不是替换系数后的Sylvester方程：左边是恒等作用减去双侧作用，向量化后的系数矩阵为$\bm I-\bm A^\top\otimes\bm A^\top$。相应地，唯一可解的条件由特征值的乘积而不是和决定。下面的定理对两侧系数不同的一般形式陈述，后面处理Schur形的局部块时要用到这一形式。

\begin{theorem}[离散型方程的唯一可解性]
\label{thm:ls-discrete-lyapunov-unique}
设$\bm G\in\mathbb R^{m\times m}$、$\bm H\in\mathbb R^{d\times d}$。方程$\bm Y-\bm G^\top\bm Y\bm H=\bm E$对每个$\bm E\in\mathbb R^{m\times d}$都有唯一解，当且仅当对$\bm G$的任一特征值$\lambda$与$\bm H$的任一特征值$\mu$，都有$\lambda\mu\neq1$。特别地，方程~\eqref{eq:ls-discrete-lyapunov}对每个$\bm Q$唯一可解，当且仅当$\bm A$在复数域中的任意两个特征值$\lambda,\lambda'$（可以是同一个）满足$\lambda\lambda'\neq1$；此时若$\bm Q$对称，则解对称。
\end{theorem}

\begin{proof}
与定理~\ref{thm:ls-sylvester-unique-solvability}的证明一样，向量化后的系数矩阵是实矩阵，唯一可解等价于它的行列式非零，所以只需在复数域中判断齐次方程$\bm Y=\bm G^\top\bm Y\bm H$是否只有零解。

必要性。设$\lambda\mu=1$。$\bm G^\top$与$\bm G$、$\bm H^\top$与$\bm H$的特征值分别相同，可取非零向量$\bm w,\bm v$使$\bm G^\top\bm w=\lambda\bm w$、$\bm H^\top\bm v=\mu\bm v$。令$\bm Y=\bm w\bm v^\top\neq\bm0$，则$\bm Y\bm H=\bm w(\bm H^\top\bm v)^\top=\mu\bm Y$，于是$\bm G^\top\bm Y\bm H=\lambda\mu\bm Y=\bm Y$，齐次方程有非零解。

充分性要用复数域中的三角化。把定理~\ref{thm:ls-real-schur}证明中的归纳搬到复数域：复方阵总有复特征向量，每一步取它张成的一维不变子空间，用标准正交基扩充成酉矩阵（第~\ref{sec:ls-adjoints}节），就得到$\bm G=\bm W_1\bm R_1\bm W_1^*$、$\bm H=\bm W_2\bm R_2\bm W_2^*$，其中$\bm W_1,\bm W_2$为酉矩阵，$\bm R_1,\bm R_2$上三角，对角元分别是$\bm G$、$\bm H$的特征值。$\bm G$是实矩阵，$\bm G^\top=\bm G^*=\bm W_1\bm R_1^*\bm W_1^*$。

设$\bm Y=\bm G^\top\bm Y\bm H$，令$\bm Z=\bm W_1^*\bm Y\bm W_2$，则$\bm Z=\bm R_1^*\bm Z\bm R_2$。$\bm R_1^*$是下三角矩阵，其$(i,k)$元为$\overline{(\bm R_1)_{k,i}}$，所以比较$(i,j)$元得
\[
 \bigl(1-\overline{(\bm R_1)_{i,i}}\,(\bm R_2)_{j,j}\bigr)z_{i,j}
 =\sum_{\substack{k\leq i,\ \ell\leq j\\(k,\ell)\neq(i,j)}}\overline{(\bm R_1)_{k,i}}\,z_{k,\ell}\,(\bm R_2)_{\ell,j}.
\]
实矩阵的非实特征值成共轭对出现，所以$\overline{(\bm R_1)_{i,i}}$仍是$\bm G$的特征值，左边系数按假设非零。右端各项满足$k+\ell<i+j$。对$i+j$归纳：$i+j=2$时右端为空，$z_{1,1}=0$；若下标和较小的元素都为零，右端为零，于是$z_{i,j}=0$。所以$\bm Z=\bm0$，$\bm Y=\bm0$。

取$\bm G=\bm H=\bm A$即得关于方程~\eqref{eq:ls-discrete-lyapunov}的结论。若$\bm Q$对称，对方程转置，由$(\bm A^\top\bm X\bm A)^\top=\bm A^\top\bm X^\top\bm A$知$\bm X^\top$也是解，由唯一性$\bm X^\top=\bm X$。
\end{proof}

下面最常用的是$\rho(\bm A)<1$的情形。此时任意两个特征值满足$|\lambda\lambda'|<1$，条件自然成立；唯一性还可以不经三角化，直接由矩阵幂的衰减证明。后面的级数论证用的就是这一直接证明。

\begin{lemma}[谱半径小于一时的唯一性]
\label{lem:ls-discrete-lyapunov-uniqueness}
若谱半径$\rho(\bm A)<1$，则方程~\eqref{eq:ls-discrete-lyapunov}至多有一个解；若$\bm Q$对称且解存在，则解对称。
\end{lemma}

\begin{proof}
两个解之差$\bm X$满足$\bm X=\bm A^\top\bm X\bm A$。反复代入得$\bm X=(\bm A^\top)^k\bm X\bm A^k$，于是$\lVert\bm X\rVert_2\leq\lVert\bm A^k\rVert_2^2\lVert\bm X\rVert_2$；由定理~\ref{thm:ls-power-decay}，$\bm A^k\to\bm0$，故$\bm X=\bm0$。若$\bm Q$对称，对方程转置可知$\bm X^\top$也是解，由唯一性$\bm X^\top=\bm X$。
\end{proof}

级数$\sum_k(\bm A^\top)^k\bm Q\bm A^k$的收敛与正定性在第~\ref{sec:ls-lyapunov-positivity}节建立。

连续型方程对应状态随时间连续变化的情形。若状态按$\dot{\bm x}=\bm A\bm x$演化（这类方程的解理论在后续卷中建立），由乘积的求导法则，二次量$\bm x^\top\bm X\bm x$的变化率为$\bm x^\top(\bm A^\top\bm X+\bm X\bm A)\bm x$。要求它恰好等于$-\bm x^\top\bm Q\bm x$，就得到连续型Lyapunov方程
\begin{equation}
 \bm A^\top\bm X+\bm X\bm A=-\bm Q .
 \label{eq:ls-continuous-lyapunov}
\end{equation}
这一段只用来说明方程的来历，本章不使用连续演化的任何结论。在代数上，它是Sylvester方程取左系数$\bm A^\top$、右系数$\bm A$、右端$-\bm Q$的特例。$\bm A^\top$与$\bm A$的特征值相同，所以Sylvester方程的谱条件在这里变成$\bm A$自身特征值两两之和的条件。

\begin{theorem}[连续型Lyapunov方程的唯一可解性]
\label{thm:ls-continuous-lyapunov-unique}
设$\bm A\in\mathbb R^{d\times d}$。方程~\eqref{eq:ls-continuous-lyapunov}对每个$\bm Q\in\mathbb R^{d\times d}$都有唯一解，当且仅当对$\bm A$在复数域中的任意两个特征值$\lambda,\lambda'$（可以是同一个），都有$\lambda+\lambda'\neq0$。此时若$\bm Q$对称，则解$\bm X$也对称。
\end{theorem}

\begin{proof}
在定理~\ref{thm:ls-sylvester-unique-solvability}中取左系数$\bm A^\top$、右系数$\bm A$。$\operatorname{spec}(\bm A^\top)=\operatorname{spec}(\bm A)$，$\operatorname{spec}(-\bm A)=\{-\lambda:\lambda\in\operatorname{spec}(\bm A)\}$，所以条件~\eqref{eq:ls-sylvester-spectral-separation}恰好是任意两个特征值之和不为零，由该定理得第一部分。若$\bm Q$对称，对方程转置得$\bm X^\top\bm A+\bm A^\top\bm X^\top=-\bm Q$，$\bm X^\top$也是解，由唯一性$\bm X^\top=\bm X$。
\end{proof}

$\bm A$的特征值实部全负时，任意两个之和的实部为负，唯一性成立。但这只是充分条件：例如$\bm A=\operatorname{diag}(1,2)$的谱和$2,3,4$都不为零，方程同样唯一可解。离散型的对应情形是：$\rho(\bm A)<1$是$\lambda\lambda'\neq1$的充分条件而非必要条件，$\bm A=\operatorname{diag}(2,3)$的各个乘积$4,6,9$都不等于$1$。两种条件的强化形式（实部全负、谱半径小于一）在第~\ref{sec:ls-lyapunov-positivity}节决定解是否正定。

\subsection{对称解的计算与检验}
\label{sec:ls-lyapunov-solve}

两类方程都可以看成Sylvester型方程，直接套用第~\ref{sec:ls-sylvester-schur}节的算法也能求解。但那样做要对$\bm A$与$\bm A^\top$分别分解，也没有利用解的对称性。本节只作一次Schur分解，并只求解的上三角部分。

\paragraph{主要思想：一次Schur变换}
对$\bm A$作实Schur分解$\bm A=\bm U\bm T\bm U^\top$，则$\bm A^\top=\bm U\bm T^\top\bm U^\top$。令$\bm Y=\bm U^\top\bm X\bm U$、$\bm D=\bm U^\top\bm Q\bm U$，$\bm D$仍对称。把分解代入方程~\eqref{eq:ls-continuous-lyapunov}与~\eqref{eq:ls-discrete-lyapunov}，左乘$\bm U^\top$、右乘$\bm U$，分别得
\[
 \bm T^\top\bm Y+\bm Y\bm T=-\bm D,\qquad
 \bm Y-\bm T^\top\bm Y\bm T=\bm D .
\]
$\bm X$对称当且仅当$\bm Y$对称，唯一可解时$\bm Y$确实对称，所以只需求$i\leq j$的元素$y_{i,j}$，下三角部分由$y_{j,i}=y_{i,j}$填入。下面先设$\bm T$上三角。

\paragraph{主要思想：连续型的逐元素递推}
$\bm T^\top$是下三角矩阵，$(\bm T^\top\bm Y)_{i,j}=\sum_{k\leq i}t_{k,i}y_{k,j}$；$(\bm Y\bm T)_{i,j}=\sum_{k\leq j}y_{i,k}t_{k,j}$。把含$y_{i,j}$的两项留在左边，得
\begin{equation}
 (t_{i,i}+t_{j,j})\,y_{i,j}
 =-d_{i,j}-\sum_{k<i}t_{k,i}y_{k,j}-\sum_{k<j}y_{i,k}t_{k,j}.
 \label{eq:ls-continuous-lyapunov-entry}
\end{equation}
按列$j=1,\ldots,d$从左向右，每列内按行$i=1,\ldots,j$从上到对角元计算。检查右端用到的元素：$y_{k,j}$（$k<i$）在同一列更靠上，已经求出；$y_{i,k}$（$k<j$）按对称性改写为上三角位置后，所在列号不超过$j$，列号等于$j$时只可能是$i=j$、行号$k<j$，也已求出。除数$t_{i,i}+t_{j,j}$是$\bm A$的两个特征值之和，由定理~\ref{thm:ls-continuous-lyapunov-unique}的条件不为零。

这一次序不能照搬第~\ref{sec:ls-sylvester-schur}节。那里左系数$\bm T$是上三角，依赖指向同列更靠下的元素，每列从下向上；这里左系数换成下三角的$\bm T^\top$，依赖指向同列更靠上的元素，所以每列从上向下。

\paragraph{主要思想：离散型的逐元素递推}
离散型中
\[
 (\bm T^\top\bm Y\bm T)_{i,j}=\sum_{k\leq i,\ \ell\leq j}t_{k,i}\,y_{k,\ell}\,t_{\ell,j},
\]
其中只有$(k,\ell)=(i,j)$一项含$y_{i,j}$。把它移到左边，得
\begin{equation}
 (1-t_{i,i}t_{j,j})\,y_{i,j}
 =d_{i,j}+\sum_{\substack{k\leq i,\ \ell\leq j\\(k,\ell)\neq(i,j)}}t_{k,i}\,y_{k,\ell}\,t_{\ell,j}.
 \label{eq:ls-discrete-lyapunov-entry}
\end{equation}
采用与连续型相同的次序。右端的$y_{k,\ell}$满足$k\leq i\leq j$、$\ell\leq j$；按对称性写成上三角位置后，所在列号$\max(k,\ell)\leq j$。列号小于$j$的元素在前面的列中已经求出；列号等于$j$时，要么$\ell=j$、$k<i$，要么$k=i=j$、$\ell<j$，两种情形都是当前列更靠上的元素。所以同一列序只读取已求元素。除数$1-t_{i,i}t_{j,j}$由定理~\ref{thm:ls-discrete-lyapunov-unique}的条件不为零。

\paragraph{主要思想：整列形式与对称性}
若不利用对称性、把整列$\bm y_j$作为未知量，两类方程按列分别写成
\[
 (\bm T^\top+t_{j,j}\bm I)\bm y_j=-\bm d_j-\sum_{\ell<j}t_{\ell,j}\bm y_\ell,\qquad
 (\bm I-t_{j,j}\bm T^\top)\bm y_j=\bm d_j+\bm T^\top\sum_{\ell<j}t_{\ell,j}\bm y_\ell .
\]
右端只用到更靠左的列，系数是下三角矩阵，每列从上到下前代。这样要求全部$d^2$个元素。若只求上三角部分，右端中左边各列的下三角元素$y_{i,\ell}$（$i>\ell$）就要取为$y_{\ell,i}$，其中$i=j$的元素恰好位于当前列，是正在求解的未知量。所以利用对称性时应按式~\eqref{eq:ls-continuous-lyapunov-entry}与~\eqref{eq:ls-discrete-lyapunov-entry}逐元素推进，未知量从$d^2$个减为$d(d+1)/2$个。

\paragraph{主要思想：二阶块与局部算子}
$\bm T$含$2\times2$对角块时，把相应的行和列成对处理。记$\bm T$的对角块为$\bm T_{I,I}$，按块写出的两类局部方程为
\[
 \begin{aligned}
 \bm T_{I,I}^\top\bm Y_{I,J}+\bm Y_{I,J}\bm T_{J,J}
 &=-\bm D_{I,J}-\sum_{K<I}\bm T_{K,I}^\top\bm Y_{K,J}-\sum_{K<J}\bm Y_{I,K}\bm T_{K,J},\\
 \bm Y_{I,J}-\bm T_{I,I}^\top\bm Y_{I,J}\bm T_{J,J}
 &=\bm D_{I,J}+\sum_{\substack{K\leq I,\ L\leq J\\(K,L)\neq(I,J)}}\bm T_{K,I}^\top\bm Y_{K,L}\bm T_{L,J}.
 \end{aligned}
\]
每个局部方程至多含四个未知数，向量化后直接消元。对角块$I=J$时$\bm Y_{I,I}$对称，只需存三个未知量。局部算子的可逆性分别由定理~\ref{thm:ls-sylvester-unique-solvability}（左系数$\bm T_{I,I}^\top$、右系数$\bm T_{J,J}$）与定理~\ref{thm:ls-discrete-lyapunov-unique}（$\bm G=\bm T_{I,I}$、$\bm H=\bm T_{J,J}$）给出：两块的特征值都是$\bm A$的特征值，所以标量情形的“分母非零”推广为“局部算子可逆”。与Sylvester方程一样，局部系数接近奇异时实现中应保留缩放因子并给出条件警告。

\paragraph{算法实现}
算法依次执行以下步骤。
\begin{enumerate}
 \item 检查$\bm Q$是否对称；不对称的输入不属于本节的问题。
 \item 计算$\bm A=\bm U\bm T\bm U^\top$与$\bm D=\bm U^\top\bm Q\bm U$。
 \item 按块列从左到右、每块列从上到对角块的次序，累计已知项，解局部方程，并把$\bm Y_{I,J}^\top$镜像填入$\bm Y_{J,I}$。
 \item 恢复$\bm X=\bm U\bm Y\bm U^\top$，在原坐标中检查残差、对称性，以及问题要求的正定性。
\end{enumerate}
第四步必须在原方程中进行，只检查变换后的小系统是不够的。也不能把算出的非对称误差强行平均后就宣布成功，对称化与残差检验回答的是不同的问题（见本节末的检验）。

\begin{example}[连续型方程的逐元素递推]
\label{ex:ls-continuous-lyapunov-entries}
取$\bm A=\bigl[\begin{smallmatrix}-1&1\\0&-2\end{smallmatrix}\bigr]$、$\bm Q=\bm I$。$\bm A$已是上三角，取$\bm U=\bm I$、$\bm T=\bm A$、$\bm D=\bm I$；特征值$-1,-2$两两之和$-2,-3,-4$都不为零。记$\bm X=\bigl[\begin{smallmatrix}p&q\\q&r\end{smallmatrix}\bigr]$，按式~\eqref{eq:ls-continuous-lyapunov-entry}依次计算：
\[
 \begin{aligned}
 (1,1):&\quad -2p=-1, & p&=\tfrac12,\\
 (1,2):&\quad -3q=-p\,t_{1,2}=-p, & q&=\tfrac16,\\
 (2,2):&\quad -4r=-1-t_{1,2}q-q\,t_{1,2}=-1-2q, & r&=\tfrac13 .
 \end{aligned}
\]
第二步就是$p-3q=0$，第三步就是$2q-4r=-1$。代回原方程：$\bm A^\top\bm X+\bm X\bm A=\bigl[\begin{smallmatrix}-2p&p-3q\\p-3q&2q-4r\end{smallmatrix}\bigr]=-\bm I$。解对称，且$p=\tfrac12>0$、$\det\bm X=\tfrac16-\tfrac1{36}=\tfrac5{36}>0$，所以$\bm X\succ0$。
\end{example}

\begin{example}[离散型方程的逐元素递推]
\label{ex:ls-discrete-lyapunov-entries}
取$\bm A=\bigl[\begin{smallmatrix}1/2&1/4\\0&1/2\end{smallmatrix}\bigr]$、$\bm Q=\bm I$，同样取$\bm U=\bm I$、$\bm T=\bm A$。特征值都是$1/2$，乘积$1/4\neq1$，分母$1-t_{i,i}t_{j,j}=3/4$。记$\bm X=\bigl[\begin{smallmatrix}p&q\\q&r\end{smallmatrix}\bigr]$，按式~\eqref{eq:ls-discrete-lyapunov-entry}依次计算。$(1,1)$元的右端没有其他项，$\tfrac34p=1$，$p=\tfrac43$。$(1,2)$元的右端只有$(k,\ell)=(1,1)$一项$t_{1,1}\,p\,t_{1,2}=p/8$，于是$\tfrac34q=p/8$，$q=\tfrac29$。$(2,2)$元的右端有三项：
\[
 t_{1,2}\,p\,t_{1,2}=\frac{p}{16},\qquad
 t_{1,2}\,q\,t_{2,2}=\frac q8,\qquad
 t_{2,2}\,q\,t_{1,2}=\frac q8,
\]
于是$\tfrac34r=1+\tfrac{p}{16}+\tfrac q4=\tfrac{41}{36}$，$r=\tfrac{41}{27}$。

代回全部元素：$\bm A^\top\bm X\bm A$的$(1,1)$元为$p/4=\tfrac13$，$(1,2)$元为$p/8+q/4=\tfrac16+\tfrac1{18}=\tfrac29$，$(2,2)$元为$p/16+q/4+r/4=\tfrac1{12}+\tfrac1{18}+\tfrac{41}{108}=\tfrac{14}{27}$。于是$\bm X-\bm A^\top\bm X\bm A$的三个元素为$\tfrac43-\tfrac13=1$、$\tfrac29-\tfrac29=0$、$\tfrac{41}{27}-\tfrac{14}{27}=1$，等于$\bm I$。又$p>0$、$\det\bm X=\tfrac43\cdot\tfrac{41}{27}-\tfrac4{81}=\tfrac{160}{81}>0$，所以$\bm X\succ0$。
\end{example}

下面的例子中$\bm A$有零特征值，还可以用累计量直接核对结果。

\begin{example}[离散型方程的递推与累计量]
\label{ex:ls-discrete-lyapunov-recursion}
取上三角矩阵$\bm A=\bigl[\begin{smallmatrix}1/2&1\\0&0\end{smallmatrix}\bigr]$与$\bm Q=\bm I$，此时可取$\bm U=\bm I$、$\bm T=\bm A$、$\bm D=\bm I$。特征值为$1/2,0$，两两乘积$1/4,0,0$都不等于$1$。按式~\eqref{eq:ls-discrete-lyapunov-entry}：$(1,1)$元满足$(1-\tfrac14)y_{1,1}=1$，得$y_{1,1}=\tfrac43$；$(1,2)$元的分母为$1-t_{1,1}t_{2,2}=1$，右端为$d_{1,2}+t_{1,1}y_{1,1}t_{1,2}=0+\tfrac12\cdot\tfrac43\cdot1$，得$y_{1,2}=\tfrac23$；$(2,2)$元的分母为$1$，右端三项中含$t_{2,2}=0$的两项为零，只剩$t_{1,2}y_{1,1}t_{1,2}=\tfrac43$，得$y_{2,2}=1+\tfrac43=\tfrac73$。于是
\[
 \bm X=\begin{bmatrix}4/3&2/3\\2/3&7/3\end{bmatrix},\qquad
 \bm A^\top\bm X\bm A=\begin{bmatrix}1/3&2/3\\2/3&4/3\end{bmatrix},
\]
两者之差为$\bm I$，残差为零。这里$\rho(\bm A)=1/2<1$，也可以从累计量核对：$\bm x_0=(0,1)^\top$时$\bm x_1=(1,0)^\top$，此后每步减半，$\sum_k\lVert\bm x_k\rVert_2^2=1+1+\tfrac14+\tfrac1{16}+\cdots=\tfrac73=x_{2,2}$。
\end{example}

前三个例子的$\bm A$已是上三角。下面的例子需要先作Schur变换。

\begin{example}[连续型方程的Schur求解]
\label{ex:ls-continuous-lyapunov-schur}
取
\[
 \bm A=\begin{bmatrix}0&3\\-1&-4\end{bmatrix},\qquad
 \bm Q=\begin{bmatrix}2&0\\0&10\end{bmatrix}.
\]
$\bm A$的特征值为$-3,-1$，谱和$-6,-4,-2$都不为零。由$\bm A(1,-1)^\top=-3(1,-1)^\top$，取$\bm U=\tfrac1{\sqrt2}\bigl[\begin{smallmatrix}1&1\\-1&1\end{smallmatrix}\bigr]$，得
\[
 \bm T=\bm U^\top\bm A\bm U=\begin{bmatrix}-3&4\\0&-1\end{bmatrix},\qquad
 \bm D=\bm U^\top\bm Q\bm U=\begin{bmatrix}6&-4\\-4&6\end{bmatrix}.
\]
按式~\eqref{eq:ls-continuous-lyapunov-entry}依次计算：$(1,1)$元$-6y_{1,1}=-6$，得$y_{1,1}=1$；$(1,2)$元$-4y_{1,2}=-d_{1,2}-y_{1,1}t_{1,2}=4-4$，得$y_{1,2}=0$；$(2,2)$元$-2y_{2,2}=-6-t_{1,2}y_{1,2}-y_{2,1}t_{1,2}=-6$，得$y_{2,2}=3$。于是$\bm Y=\operatorname{diag}(1,3)$，变回原坐标得
\[
 \bm X=\bm U\bm Y\bm U^\top=\begin{bmatrix}2&1\\1&2\end{bmatrix}.
\]
核验：$\bm A^\top\bm X=\bigl[\begin{smallmatrix}-1&-2\\2&-5\end{smallmatrix}\bigr]$，加上它的转置$\bm X\bm A$得$\operatorname{diag}(-2,-10)=-\bm Q$，原方程残差为零，$\bm X$对称。
\end{example}

\paragraph{算法分析：对称性、残差与正定性检验}
对需要正定解的应用，还要另外检查对称性与正定性，并在浮点计算中把三项检验放在一起解释。算出的$\bm X$可能有舍入量级的微小不对称。对称化只消除这一偏差，不能说明方程解对了；后者要看原方程的相对残差是否在机器精度的若干倍以内。

正定性可以用Cholesky分解（第~\ref{sec:ls-cholesky}节）检验，但浮点中分解成功只说明一个与$\bm X$相差舍入量级的对称矩阵正定（第~\ref{sec:ls-numerical-stability}节）。只有当$\lambda_{\min}(\bm X)$明显大于这一舍入量级时，分解成败才能判定$\bm X$本身是否正定；否则还应结合残差容差和问题本身是否保证正定（第~\ref{sec:ls-lyapunov-positivity}节）来判断。示例~\ref{ex:ls-continuous-lyapunov-schur}中$\bm X$的Cholesky因子对角元为$\sqrt2$与$\sqrt{3/2}$，分解成功，$\bm X\succ0$。

\paragraph{算法分析：计算代价}
一次实Schur分解约需$25d^3$次运算，变换右端与恢复解各约$O(d^3)$次。连续型递推中每个元素的和式有$O(d)$项，总量$O(d^3)$。离散型的式~\eqref{eq:ls-discrete-lyapunov-entry}含双重和，逐元素直接求和总量为$O(d^4)$；实现时按列保存部分乘积（如$\sum_{\ell<j}t_{\ell,j}\bm y_\ell$及其左乘$\bm T^\top$的结果），可以把总量降到$O(d^3)$。所以两类方程的总代价都由$O(d^3)$主导，存储为$O(d^2)$。

\subsection{正定解的存在条件}
\label{sec:ls-lyapunov-positivity}

上一节的算法在唯一可解时算出对称解，但不保证解正定。本节先回到离散累计：$\rho(\bm A)<1$时，离散型方程的解就是第~\ref{sec:ls-lyapunov-solvability}节累计量$\bm X_N$的极限，正定性可以从级数逐项读出。

\begin{theorem}[离散型Lyapunov方程的级数解]
\label{thm:ls-discrete-lyapunov-series}
设$\rho(\bm A)<1$，$\bm Q$对称。则级数
\begin{equation}
 \bm X=\sum_{k=0}^{\infty}(\bm A^\top)^k\bm Q\bm A^k
 \label{eq:ls-discrete-lyapunov-series}
\end{equation}
收敛，且是方程~\eqref{eq:ls-discrete-lyapunov}的唯一解。若$\bm Q\succeq0$，则$\bm X\succeq0$；若$\bm Q\succ0$，则$\bm X\succ0$。
\end{theorem}

\begin{proof}
取$r\in(\rho(\bm A),1)$。由引理~\ref{lem:ls-power-geometric-bound}，存在$C_r$使$\lVert\bm A^k\rVert_2\leq C_rr^k$，所以第$k$项的范数不超过$\lVert\bm Q\rVert_2C_r^2r^{2k}$。这是收敛几何级数的通项，由定理~\ref{thm:ls-absolute-series-convergence}，级数绝对收敛，记部分和为$\bm X_N$。

映射$\mathcal M(\bm Y)=\bm A^\top\bm Y\bm A$满足$\lVert\mathcal M(\bm Y)\rVert_2\leq\lVert\bm A\rVert_2^2\lVert\bm Y\rVert_2$，是连续的线性映射，故$\mathcal M(\bm X_N)\to\mathcal M(\bm X)$。逐项作用后移项相消，
\[
 \bm X_N-\mathcal M(\bm X_N)=\bm Q-(\bm A^\top)^{N+1}\bm Q\bm A^{N+1},
\]
右端末项的范数不超过$\lVert\bm Q\rVert_2C_r^2r^{2(N+1)}\to0$。令$N\to\infty$得$\bm X-\bm A^\top\bm X\bm A=\bm Q$，唯一性由引理~\ref{lem:ls-discrete-lyapunov-uniqueness}给出。

对任意$\bm x$，$\bm x^\top(\bm A^\top)^k\bm Q\bm A^k\bm x=(\bm A^k\bm x)^\top\bm Q(\bm A^k\bm x)$，所以$\bm Q\succeq0$时各项都半正定，$\bm X\succeq0$。$\bm Q\succ0$时，对$\bm x\neq\bm0$，首项$\bm x^\top\bm Q\bm x>0$而其余项非负，故$\bm X\succ0$。
\end{proof}

标量情形可以直接核对符号和分母。$d=1$、$|a|<1$时，级数为$x=\sum_kq\,a^{2k}=q/(1-a^2)$，代入方程得$x-a^2x=q$；分母$1-a^2$正是唯一性条件中的$1-\lambda\lambda$，$q>0$时$x>0$。

\paragraph{半正定右端}
实际问题中常只计量状态的一部分，例如$\bm Q=\bm c\bm c^\top$只记录一个线性读数$\bm c^\top\bm x_k$的平方，这时$\bm Q$奇异。定理~\ref{thm:ls-discrete-lyapunov-series}只保证$\bm X\succeq0$。$\bm X$是否正定，取决于是否存在一个非零初始方向，沿它出发的整条轨迹始终不被累计量察觉。

\begin{theorem}[半正定右端的正定条件]
\label{thm:ls-discrete-lyapunov-semidefinite}
设$\rho(\bm A)<1$，$\bm Q\succeq0$，$\bm X$为级数~\eqref{eq:ls-discrete-lyapunov-series}给出的解。对$\bm x\in\mathbb R^d$，$\bm x^\top\bm X\bm x=0$当且仅当$\bm Q\bm A^k\bm x=\bm0$对$k=0,1,\ldots,d-1$都成立。因此$\bm X\succ0$，当且仅当
\[
 \begin{bmatrix}\bm Q\\\bm Q\bm A\\\vdots\\\bm Q\bm A^{d-1}\end{bmatrix}
\]
的零空间只含零向量，即它的秩为$d$。
\end{theorem}

\begin{proof}
由级数，$\bm x^\top\bm X\bm x=\sum_{k\geq0}(\bm A^k\bm x)^\top\bm Q(\bm A^k\bm x)$，各项非负，所以它为零当且仅当每一项为零。对半正定的$\bm Q$，$\bm v^\top\bm Q\bm v=0$等价于$\bm Q\bm v=\bm0$：写$\bm Q=\bm L\bm L^\top$（例如取谱分解的平方根），$\bm v^\top\bm Q\bm v=\lVert\bm L^\top\bm v\rVert_2^2=0$推出$\bm L^\top\bm v=\bm0$，从而$\bm Q\bm v=\bm0$；反向显然。于是$\bm x^\top\bm X\bm x=0$当且仅当$\bm Q\bm A^k\bm x=\bm0$对所有$k\geq0$成立。

只需检查前$d$个$k$：由Cayley--Hamilton定理~\ref{thm:ls-cayley-hamilton}，$\bm A^d$是$\bm I,\bm A,\ldots,\bm A^{d-1}$的线性组合；乘以$\bm A$后归纳可知，每个$\bm A^k$（$k\geq d$）都是这些幂的线性组合。所以$\bm Q\bm A^k\bm x=\bm0$对$k<d$成立时，对所有$k$都成立。最后，$\bm X\succeq0$，所以$\bm X\succ0$当且仅当$\bm x^\top\bm X\bm x=0$只有零解，即上述堆叠矩阵的零空间只含零向量。
\end{proof}

这个条件在控制理论中称为可观测性：从读数序列$\bm Q\bm A^k\bm x$能否区分初始状态。它只依赖$\bm A$与$\bm Q$的零空间结构，不依赖$\bm Q$的具体大小。

\begin{example}[未被累计量察觉的方向]
\label{ex:ls-discrete-lyapunov-unobserved}
取$\bm A=\bigl[\begin{smallmatrix}0&1\\0&0\end{smallmatrix}\bigr]$，$\rho(\bm A)=0$，且$\bm A^2=\bm0$，级数只有两项$\bm X=\bm Q+\bm A^\top\bm Q\bm A$。$\bm A$把$\bm e_2$送到$\bm e_1$，把$\bm e_1$送到零。

若$\bm Q=\operatorname{diag}(1,0)$，只计量第一个坐标，则$\bm A^\top\bm Q\bm A=\operatorname{diag}(0,1)$，$\bm X=\bm I\succ0$。从$\bm e_2$出发时第一步没有被计量，但下一步状态移到$\bm e_1$上被察觉；按定理~\ref{thm:ls-discrete-lyapunov-semidefinite}，$\bm Q\bm x=\bm0$给出$x_1=0$，$\bm Q\bm A\bm x=(x_2,0)^\top=\bm0$又给出$x_2=0$。

若$\bm Q=\operatorname{diag}(0,1)$，只计量第二个坐标，则$\bm Q\bm A=\bm0$，$\bm X=\bm Q=\operatorname{diag}(0,1)$奇异。方向$\bm e_1$满足$\bm Q\bm e_1=\bm0$、$\bm Q\bm A\bm e_1=\bm0$：从它出发的轨迹$\bm e_1,\bm0,\bm0,\ldots$始终不被计量，$\bm e_1^\top\bm X\bm e_1=0$。两种右端都只是秩一的半正定矩阵，结果却不同，区别在于未计量的方向是否会被$\bm A$送进计量范围。
\end{example}

\paragraph{正定解给出的谱证书}
前面由谱条件推出正定解。反过来，一个对称正定解本身就证明了谱条件，而且证明只用代数运算。这给上一节算法的输出提供了可检验的证书：残差小、Cholesky分解在可靠的容差下成功，就同时检验了$\bm A$的稳定性。

\begin{theorem}[正定解给出的谱证书]
\label{thm:ls-lyapunov-spectral-certificate}
设$\bm Q\succ0$，$\bm X$对称正定。
\begin{enumerate}
 \item 若$\bm X-\bm A^\top\bm X\bm A=\bm Q$，则$\bm A$的每个特征值满足$|\lambda|<1$，即$\rho(\bm A)<1$。
 \item 若$\bm A^\top\bm X+\bm X\bm A=-\bm Q$，则$\bm A$的每个特征值的实部为负。
\end{enumerate}
\end{theorem}

\begin{proof}
设$\bm A\bm v=\lambda\bm v$，$\bm v=\bm a+\mathrm i\bm b\neq\bm0$，$\bm a,\bm b$为实向量。对实对称矩阵$\bm M$，$\bm v^*\bm M\bm v=\bm a^\top\bm M\bm a+\bm b^\top\bm M\bm b$，交叉项$\mathrm i(\bm a^\top\bm M\bm b-\bm b^\top\bm M\bm a)$因对称而抵消；所以$\bm v^*\bm X\bm v>0$、$\bm v^*\bm Q\bm v>0$。$\bm A$是实矩阵，$\bm v^*\bm A^\top=(\bm A\bm v)^*=\bar\lambda\bm v^*$。

离散型：$\bm v^*\bm A^\top\bm X\bm A\bm v=\bar\lambda\lambda\,\bm v^*\bm X\bm v$，于是$(1-|\lambda|^2)\,\bm v^*\bm X\bm v=\bm v^*\bm Q\bm v>0$，得$|\lambda|<1$。

连续型：$\bm v^*(\bm A^\top\bm X+\bm X\bm A)\bm v=(\bar\lambda+\lambda)\,\bm v^*\bm X\bm v=2\operatorname{Re}\lambda\,\bm v^*\bm X\bm v$，它等于$-\bm v^*\bm Q\bm v<0$，得$\operatorname{Re}\lambda<0$。
\end{proof}

示例~\ref{ex:ls-continuous-lyapunov-entries}与~\ref{ex:ls-discrete-lyapunov-entries}可以按这个证书反向读：两个$\bm X$都对称正定、$\bm Q=\bm I$，分别证明了$\bm A=\bigl[\begin{smallmatrix}-1&1\\0&-2\end{smallmatrix}\bigr]$的特征值实部为负、$\bm A=\bigl[\begin{smallmatrix}1/2&1/4\\0&1/2\end{smallmatrix}\bigr]$的谱半径小于一；对这两个三角矩阵，这与直接读对角元的结论一致。

\paragraph{连续型的正定条件}
连续型方程有与定理~\ref{thm:ls-discrete-lyapunov-series}平行的结论：当$\bm A$的特征值实部全负且$\bm Q\succ0$时，唯一解$\bm X$正定；$\bm Q\succeq0$时$\bm X\succeq0$，并且与定理~\ref{thm:ls-discrete-lyapunov-semidefinite}同样，$\bm X\succ0$当且仅当$\bm Q\bm A^k\bm x=\bm0$（$k=0,\ldots,d-1$）只有零解。这里的唯一性由定理~\ref{thm:ls-continuous-lyapunov-unique}保证，正定性则需要另外的论证：把离散级数换成积分$\bm X=\int_0^\infty e^{t\bm A^\top}\bm Qe^{t\bm A}\,\mathrm dt$，再验证积分收敛并满足方程。

这一论证要用矩阵指数$e^{t\bm A}$的定义、导数以及特征值实部全负时的指数衰减，可参考\BookChapterRef{chap:linear-differential-systems}{卷六“线性微分方程与系统”章}；积分表示与轨迹长期稳定性的完整推导在\BookChapterRef{chap:dynamical-systems-long-term}{卷九“动力系统与轨迹的长期行为”章}给出。本章只保留条件陈述和代数检验，求解算法也不以它为前提。

低维例子可以直接核验这一陈述。标量方程$2ax=-q$在$a<0$、$q>0$时给出$x=q/(2|a|)>0$。示例~\ref{ex:ls-continuous-lyapunov-schur}中$\bm A$的特征值为$-3,-1$，$\bm Q=\operatorname{diag}(2,10)\succ0$，算出的$\bm X$特征值为$1,3$，确实正定；由定理~\ref{thm:ls-lyapunov-spectral-certificate}，这个正定解反过来也证明了$\bm A$的特征值实部为负。
